\section{Methods}
\label{sec:methods}

% This is the section that earns or loses the paper. Everything here should be
% traceable to a workflow step (WF-xxx) in the planning workbook, and every
% subsection should end able to answer "how would someone reproduce this?"

\subsection{Study area and analysis extent}
\label{sec:extent}
\TODO{WF-001. Countries, area, the confirmed-range envelope, exclusion rules, and the
statement that the extent was fixed before any connectivity output was viewed.}
Analyses were conducted in \TODO{equal-area CRS, e.g. Africa Albers Equal Area Conic,
give the full WKT in Appendix~\ref{app:protocol}} at a working resolution of \cellsize.

\subsection{Rationale for the working resolution}
\label{sec:scale}
\TODO{CAU-006. The occurrence baseline is a 10-km assessment grid; environmental
covariates range from 30\,m to 1\,km. Justify \cellsize\ as the compromise: fine
enough that a pinch point is a meaningful location, coarse enough that it does not
imply precision the occurrence evidence cannot support. State that the full pipeline
was repeated at 2\,km and 5\,km as a resolution sensitivity (WF-022).}

\subsection{Range cores}
\label{sec:cores}
\TODO{WF-011. Three or more core definitions from the archived confirmed-range data,
differing in evidence-category threshold and minimum patch area. Report how many
cores each definition yields. State explicitly that cores are HELD CONSTANT across
all four temporal snapshots, and why (CAU-002): no comparable range reassessment
exists after the baseline period, so any core movement would be invented.}

\subsection{Covariates and their temporal treatment}
\label{sec:covariates}
\TODO{Table~\ref{tab:covariates} lists every layer, its native resolution, its
source year(s), and critically whether it is treated as DYNAMIC (varies across
snapshots) or STATIC (held fixed, contributing no temporal signal). Do not bury this
distinction. Roads, livestock, terrain and hydrography are static; built-up fraction,
nighttime radiance and continuous vegetation fields are dynamic.}

\paragraph{Why categorical land cover is not the change carrier.}
\TODO{CAU-008 and the MODIS class-instability problem. Annual categorical products
flip classes between years for reasons unrelated to landscape change. Temporal signal
is therefore carried on continuous fields; categorical land cover is used for masking
and description only.}

\subsection{Resistance scenarios}
\label{sec:resistance}
\TODO{WF-013, WF-014. Present RES-01 to RES-05 as a table of variable-to-resistance
functions with the evidence or rationale for each. State that the scenario set was
fixed before results were viewed and that no single scenario is designated as
preferred. Report the transformation (linear vs exponential) per scenario and cite
\textcite{zeller2012resistance} for why the transformation matters as much as the
weights.}

\subsection{Connectivity modelling}
\label{sec:connectivity}
\TODO{WF-015, WF-016. Least-cost / effective-distance work in ArcGIS Pro (Distance
Accumulation, Optimal Region Connections, Corridor). Current flow computed with
Circuitscape.jl \parencite{anantharaman2020circuitscape} in pairwise mode; give the
solver, the source/ground configuration, and how current was normalised for
comparability across years. State why the omnidirectional approach
\parencite{landau2021omniscape} was NOT used: cores are defined a priori, so a
core-to-core formulation is the correct one.}

\paragraph{Corridor definition.}
\TODO{Decide and state ONCE whether corridors are (a) continuous current surfaces
summarised by quantile or (b) thresholded polygons. This choice sets the denominator
for the protected-area comparison in \S\ref{sec:pa} and cannot be made after seeing
the overlap result.}

\subsection{Network structure}
\label{sec:graph}
\TODO{WF-017. Cores as nodes, effective distances as edge weights. Report the
probability of connectivity and per-node $dPC$ \parencite{saura2007pc}, computed in
Conefor \parencite{saura2009conefor} or an equivalent, and betweenness on the same
graph. Justify the metric set against \textcite{keeley2021metrics}. Report
sensitivity to the dispersal-distance parameter; do not report a single value.}

\subsection{Temporal comparison}
\label{sec:temporal}
\TODO{WF-018. Same cores, same scenarios, same resolution across all four snapshots,
so that differences are attributable to covariate change alone. Give the change
equations. Report both absolute and proportional change.}

\subsection{Protected-area coverage}
\label{sec:pa}
\TODO{WF-019, CAU-011, CAU-012. Name the \wdpaversion. State polygon-only handling
and what was excluded. State the null explicitly: high-current area is compared
against \TODO{the full extent OR the modelled network} using an area-standardised
permutation test. Predeclare this.}

\subsection{Robustness and uncertainty}
\label{sec:uncertainty}
\TODO{WF-022. The full sensitivity grid: \nscen\ resistance scenarios $\times$ core
definitions $\times$ resolutions $\times$ snapshots. Report the resulting run count.
Define "robust" numerically (e.g. a cell in the top decile of current under at least
$k$ of $n$ runs) and fix $k$ in advance. Produce an agreement map. Low-agreement
locations are reported as uncertain, never as priorities.}

\subsection{Independent plausibility assessment}
\label{sec:validation}
\TODO{WF-020, WF-021, CAU-004, CAU-017. See the validation decision note before
writing this. If the plausibility-check route is taken, this subsection is titled as
above and must NOT use the word "validation". If the empirical-resistance route is
taken, add RES-06, describe the fitting procedure, and report country-held-out
performance following \textcite{roberts2017cv}.}

\subsection{Monitoring prioritisation}
\label{sec:mcda}
\TODO{WF-023, WF-024, CAU-020. Components, normalisation, and the predeclared weight
set, plus at least two alternative weightings. Report rank stability (Spearman /
Kendall) across weightings. Publish component scores alongside the composite so a
reader can reweight.}

\subsection{Software and reproducibility}
\label{sec:software}
\TODO{ArcGIS Pro \TODO{version}; Circuitscape.jl \TODO{version}; Julia \TODO{version};
R \TODO{version} or Python \TODO{version}; Google Earth Engine for acquisition and
harmonisation. All geoprocessing scripted in arcpy rather than ModelBuilder. Archive
location and what can and cannot be redistributed (CAU-012, CAU-013).}
